%% korr-preamble.tex -- gemeinsame Präambel für die Korrekturlisten LNM 1184
%% Makros wie im Buch (LN-Book-Springer/preamble), damit die Formeln
%% genauso aussehen wie im Proof.
\usepackage[T1]{fontenc}
\usepackage[utf8]{inputenc}
\usepackage{newtxtext}
\usepackage[varvw]{newtxmath}
\usepackage[scaled=.9]{helvet}
\usepackage[varqu,scaled=.95]{zi4}
\usepackage[ngerman]{babel}
\usepackage{csquotes}
\usepackage[a4paper,landscape,margin=1.6cm,top=1.8cm,bottom=1.8cm]{geometry}
\usepackage{microtype}
\usepackage{array,longtable,booktabs}
\usepackage{ragged2e}
\usepackage[table]{xcolor}
\usepackage{xspace}
\usepackage{enumitem}
\setlist{nosep,leftmargin=*}
\usepackage[nointegrals]{wasysym}
\usepackage{newunicodechar}
\usepackage{nicefrac}
\usepackage{mathtools}
\usepackage{mleftright}
\mleftright
\usepackage[hidelinks]{hyperref}
\setlength{\parindent}{0pt}
\setlength{\parskip}{.5em}
\setlength{\emergencystretch}{3em}
\renewcommand{\arraystretch}{1.25}
\definecolor{codegray}{gray}{0.35}

%% ---- Makros aus dem Buch -------------------------------------------
\DeclarePairedDelimiterX{\norm}[1]{\lVert}{\rVert}{\ifblank{#1}{\:\cdot\:}{#1}}
\DeclarePairedDelimiterX{\abs}[1]{\lvert}{\rvert}{\ifblank{#1}{\:\cdot\:}{#1}}
\DeclarePairedDelimiterX{\dualp}[1]{\langle}{\rangle}{\ifblank{#1}{\,\cdot\,,\cdot\,}{#1}}
\newcommand{\interval}[1]{\left[ #1 \right]}
\newcommand{\ointerval}[1]{\left] #1 \right[}
\newcommand{\rointerval}[1]{\left[ #1 \right[}
\newcommand{\lointerval}[1]{\left] #1 \right]}
\newcommand{\tint}{\textstyle\int}
\AtBeginDocument{%
\let\ORGvarepsilon=\varepsilon
\let\epsilon=\epsilon
\let\epsilon=\ORGvarepsilon
\let\ORGvarrho=\varrho
\let\rho=\rho
\let\rho=\ORGvarrho
\let\ORGvartheta=\vartheta
\let\theta=\theta
\let\theta=\ORGvartheta
\let\ORGvarphi=\varphi
\let\phi=\phi
\let\phi=\ORGvarphi
\let\ORGvarleq=\leqslant
\let\leqslant=\leq
\let\leq=\ORGvarleq
\let\ORGvargeq=\geqslant
\let\geqslant=\geq
\let\geq=\ORGvargeq
\let\ORGvarpreccurlyeq=\preccurlyeq
\let\preccurlyeq=\preceq
\let\preceq=\ORGvarpreccurlyeq
\let\setminus\smallsetminus			% A - B fuer Mengen
\renewcommand{\P}{\mathfrak{P}}		% Achtung: \P hat andere Bedeutung in LaTeX
}
\newcommand{\N}{\mathbb{N}}		% Natuerliche Zahlen
\newcommand{\Z}{\mathbb{Z}}		% Ganze Zahlen
\newcommand{\Q}{\mathbb{Q}}		% Rationale Zahlen
\newcommand{\R}{\mathbb{R}}		% Reelle Zahlen
\newcommand{\C}{\mathbb{C}}		% Komplexe Zahlen
\newcommand{\K}{\mathbb{K}}		% Koerperzeichen
\newcommand{\T}{\mathbb{T}}		% Torus
\DeclareMathOperator{\Id}{Id}
\DeclareMathOperator{\sign}{sign}				% signum-Operator
\newcommand{\im}{\ensuremath{\mathrm{i}}}		% imaginäre Konstante i
\newcommand{\iu}{\ensuremath{\mathrm{i}}}
\newcommand{\eu}{\ensuremath{\mathrm{e}}}
\renewcommand*{\Re}{\mathop{}\!\mathrm{Re}}		% Realteil
\renewcommand*{\Im}{\mathop{}\!\mathrm{Im}}		% Imaginärteil
\newcommand*{\supp}{\mathop{}\!\mathrm{supp}\,}  % Träger
\newcommand*{\diff}[1]{\mathop{}\!\mathrm{d}{#1}}	% \diff{s} = ds, 
\newcommand*{\ds}{\mathop{}\!\mathrm{d}{s}}       	% \ds = ds, 
\newcommand*{\dg}{\mathop{}\!\mathrm{d}{g}}       	% \dg = dg, 
\newcommand*{\dt}{\mathop{}\!\mathrm{d}{t}}       	% \dt = dt, 
\newcommand*{\dx}{\mathop{}\!\mathrm{d}{x}}       	% \dx = dx, 
\newcommand*{\dy}{\mathop{}\!\mathrm{d}{y}}       	% \dy = dy, 
\newcommand*{\dr}{\mathop{}\!\mathrm{d}{r}}       	% \dr = dr, 
\newcommand*{\dm}{\mathop{}\!\mathrm{d}{m}}       	% \dm = dm, 
\newcommand*{\du}{\mathop{}\!\mathrm{d}{u}}       	% \du = du, 
\DeclareMathOperator{\esssup}{ess\,sup}
\newcommand{\LE}{\mathcal{L}(E)} 		% 
\newcommand{\BH}{\mathcal{B}(H)} 		% 
\newcommand{\BA}{\mathfrak{A}}
\newcommand{\ZE}{\mathcal{Z}(E)} 	% \ZE  Zentrum von E
\newcommand{\LH}{\mathcal{L}(H)} 
\let\originalL\L
\AtBeginDocument{%
  \DeclareRobustCommand{\L}[1]{\mathcal{L}(#1)}%
}
\newcommand{\Fix}[1]{\mathop{}\mathrm{Fix}\left(#1\right)}		% \Fix{T} Fixraum
\newcommand{\rank}{\mathop{}\!\mathrm{rank}\,}		% 
\newcommand{\Kern}[1]{\mathop{}\mathrm{ker}(#1)}
\newcommand{\Image}[1]{\mathop{}\mathrm{im}(#1)}
\newcommand{\image}{\mathop{}\!\mathrm{im}\,}
\renewcommand{\dim}{\mathop{}\!\mathrm{dim}\,}
\newcommand*{\eg}{e.g.,\xspace}
\newcommand*{\ie}{i.e.,\xspace}
\newcommand*{\resp}{resp.\xspace}
\newcommand*{\vs}{vs.\xspace}
\newcommand*{\cf}{cf.\xspace}
\newcommand*{\etc}{etc.\xspace}
\newcommand{\CA}{$\mathrm{C}^{*}$}	% C*-Algebra
\newcommand{\WA}{$\mathrm{W}^{*}$}	% W*-Algbera
\newcommand{\AW}{$\mathrm{AW}^{*}$}	% AW*-Algbera
\newcommand*{\TT}{\mathcal{T}}		% C_0-Semigroup
\newcommand*{\RR}{\mathcal{R}}		% Pseudo Resolvente
\newcommand*{\SG}{\mathcal{S}}	    % Semigroup
\newcommand*{\UG}{\mathcal{U}}	    % Implemented Semigroup
\newcommand*{\F}{\mathcal{F}}		% F-product
\newcommand{\mfrac}[2]{%
    {\raise.5ex\hbox{$\scriptstyle #1$}}%
    \kern.025em/\kern.025em%
    {\lower.25ex\hbox{$\scriptstyle #2$}}}
\newcommand*{\sfrac}[2]{\leavevmode\kern.1em%
       \raise.5ex\hbox{\small #1}%
       \kern.025em/\kern.025em%
       \lower.25ex\hbox{\small #2}}
\newcommand{\1}{\mathbf{1}}
\newcommand{\HO}{\mathcal{H}(O)} 		%  für Lotz
\newcommand{\WO}{\mathcal{W}(O)} 		%  für Lotz
\newcommand*{\ddp}{\mathop{}\!\mathrm{d}{p}}       	% \ddp = dp, 
\newcommand*{\dN}{\mathop{}\!\mathrm{d}{N}}       	% \dN = dN, 
\newcommand*{\ddt}{\mathop{}\!\mathrm\fraq{d}{dt}}   % \ddt = d/dt
\newcommand*{\lnm}{LNM1184: } % \LMN-BookCode
\newcommand{\closure}[1]{#1^{-}}
%% Hilfsmatrizen aus Teil I, Kap. 1
\newcommand{\ABC}{\begin{pmatrix} 0 & 1 \\ -1 & 0 \end{pmatrix}}
\newcommand{\BBC}{\begin{pmatrix} \cos(t) & \sin(t) \\ - \sin(t) & \cos(t) \end{pmatrix}}
\newcommand{\CBC}{\begin{pmatrix} 1 & 1 \\ -1 & -1 \end{pmatrix}}
\newcommand{\DBC}{\begin{pmatrix} 1+t & t \\ - t & 1-t \end{pmatrix}}
\newcommand{\EBC}{\begin{pmatrix} 0 & 1 \\ 1 &0 \end{pmatrix}}
\newcommand{\FBC}{\begin{pmatrix} \cosh(t) & \sinh(t) \\ \sinh(t) & \cosh(t) \end{pmatrix}}

%% ---- Darstellung der Listen ----------------------------------------
%% Quelltext, der sich nicht als Formel setzen lässt (Labels, \ref, Dateinamen)
\newcommand{\code}[1]{\leavevmode{\color{codegray}\ttfamily\small #1}}
\newcommand{\sichA}{\CIRCLE}      % ● hoch
\newcommand{\sichB}{\LEFTcircle}  % ◐ mittel
\newcommand{\sichC}{\Circle}      % ○ niedrig
%% Querverweise im Buch: Label-Namen statt Nummern
\newcommand{\lbl}[1]{\textsf{\footnotesize\color{codegray}\detokenize{#1}}}
